% Part: first-order-logic
% Chapter: proof-systems
% Section: seqeunt-calculus

\documentclass[../../../include/open-logic-section]{subfiles}

\begin{document}

\iftag{FOL}
      {\olfileid{fol}{prf}{seq}}
      {\olfileid{pl}{prf}{seq}}

\olsection{Kalkulus Sekuen}

Nalika akeh sistem !!{derivation} nggarap susunan !!{sentence}s,
kalkulus sekuen nggarap \emph{sekuen}. Sekuen yaiku ekspresi kanthi
wujud
\[
!A_1, \dots, !A_m \Sequent !B_1, \dots, !B_m,
\]
yaiku pasangan urutan !!{sentence}s kang dipisahake dening simbol
sekuen~$\Sequent$. Salah siji urutan bisa kosong. !!^a{derivation} ing
kalkulus sekuen iku wit sekuen; sekuen paling ndhuwur nduweni wujud
khusus (diarani ``sekuen wiwitan'' utawa ``aksioma''), lan saben sekuen
liyane ndherek sekuen kang persis ana ing ndhuwure nganggo salah siji
aturan inferensi. Aturan inferensi bisa ngowahi !!{sentence}s ing
sekuen (nambah, mbusak, utawa nata maneh ing sisih kiwa utawa tengen),
utawa ngenalake !!{formula} kompleks ing kesimpulan aturan. Tuladhane,
aturan $\LeftR{\land}$ ngidini inferensi saka $!A,
\Gamma \Sequent \Delta$ menyang $!A \land !B, \Gamma \Sequent \Delta$, lan
$\RightR{\lif}$ ngidini inferensi saka $!A, \Gamma \Sequent
\Delta, !B$ menyang $\Gamma \Sequent \Delta, !A \lif !B$, kanggo saben
$\Gamma$, $\Delta$, $!A$, lan~$!B$. (Mligine, $\Gamma$ lan~$\Delta$
bisa kosong.)

Relasi $\Proves$ adhedhasar kalkulus sekuen ditemtokake kaya mangkene:
$\Gamma \Proves !A$ yen lan mung yen ana sawijining urutan $\Gamma_0$
saengga saben $!A$ ing $\Gamma_0$ kalebu ing~$\Gamma$ lan ana
!!{derivation} kang oyode ngemot sekuen~$\Gamma_0 \Sequent !A$.
$!A$ iku teorema ing kalkulus sekuen yen sekuen~$\Sequent !A$ duwe
!!a{derivation}. Tuladhane, iki !!a{derivation} kang nuduhake yen
$\Proves (!A \land !B) \lif !A$:
\begin{prooftree}
  \Axiom$!A \fCenter !A$
  \RightLabel{\LeftR{\land}}
  \UnaryInf$!A \land !B \fCenter !A$
  \RightLabel{\RightR{\lif}}
  \UnaryInf$\fCenter (!A \land !B) \lif !A$
\end{prooftree}

Himpunan $\Gamma$ ora konsisten ing kalkulus sekuen yen ana
!!a{derivation} saka $\Gamma_0 \Sequent$ (saben $!A \in \Gamma_0$
kalebu ing~$\Gamma$ lan sisih tengen sekuen kosong). Kanthi nggunakake
aturan \RightR{\Weakening}, saben !!{sentence} bisa !!{derive}d saka
himpunan kang ora konsisten.

Kalkulus sekuen ditemokake ing taun 1930-an dening Gerhard Gentzen.
Amarga rancangane sistematis lan simetris, kalkulus iki dadi formalisme
kang migunani banget kanggo ngembangake teori !!{derivation}s. Golek
!!{derivation}s ing kalkulus sekuen cukup gampang, nanging
!!{derivation}s iki kerep angel diwaca lan gegayutane karo bukti
kadhangkala ora gampang katon. Nanging, kalkulus iki kabukten dadi
pendekatan kang endah banget kanggo sistem !!{derivation}, lan akeh
logika duwe sistem kalkulus sekuen.

\end{document}
